\begin{afterword}
\summaryhead

The post opens at the funeral of the author's great-grandmother. A grand-uncle said that God had taken her piece by piece. The author was surprised, having believed that Jews usually avoid attributing tragedies to God. The grand-uncle later left the religion.

The author describes Modern Orthodox Judaism as a tradition that permits doubt but only of a safe kind. Believers raise objections they can answer, such as the conflict between the six days of creation and the age of the universe. They do not dwell on harder ones, such as God killing the firstborn of Egypt, which the Passover Seder commemorates.

The post generalizes. People who question a cherished belief tend to attack it where a comforting reply is ready and to stop at the first reply. The avoidance is instinctive, like pulling back from a hot stove. Its remedy is to think deliberately about whatever hurts most, to ask what smart opponents would say to one's replies, and to recite Gendlin's lines about facing the truth.

\responsehead

The post teaches a method for questioning your own beliefs and demonstrates it only on other people's. The belief under examination is the religion of the author's family. No belief of the author's is tested. The evidence for the theory begins with a grieving son's eulogy for his mother, read as if it were a statement of doctrine.

The post's account of the religion is wrong where it matters most. Its premise is that Jews avoid attributing tragedy to God. Jewish law requires the opposite: a blessing for bad news as for good, said on hearing of a death and by mourners at the funeral. Its story about the Seder wine credits old rabbis with a sympathy that was first written down in 1904. The post claims to know how Orthodox Jews think while quoting none of them.

Turn the post's method on the post. The post rests on a belief of the author's: that educated believers stay religious by avoiding their faith's weak points. Where is that belief weakest? At the question of what Jews actually do with the hard passages. That is the question the post avoids. It grants that there are ``probably a dozen'' replies about the firstborn and reads none. It offers an imagined believer's inner monologue in place of any real one. In my reading it attacks its target only where it knows it can win, which is the post's own definition of avoidance.

The thesis also rests on almost nothing. Anyone who thinks about the slaughtered firstborn ``is \emph{really} questioning it'' and ``probably'' will soon leave. That could be tested, by asking believers who have faced the hard case whether they stayed, but the post's evidence is one relative who left.

In short: a lesson in attacking one's own beliefs, demonstrated on the author's relatives' faith, built on errors that the faith's own texts correct, and never applied to the author.
\end{afterword}
